\documentclass[10pt]{article}
\usepackage[utf8]{inputenc}
\usepackage[T1]{fontenc}
\usepackage{amsmath}
\usepackage{amsfonts}
\usepackage{amssymb}
\usepackage[version=4]{mhchem}
\usepackage{stmaryrd}
\usepackage{bbold}

\DeclareUnicodeCharacter{21D2}{\ifmmode\Rightarrow\else{$\Rightarrow$}\fi}
\DeclareUnicodeCharacter{2192}{\ifmmode\rightarrow\else{$\rightarrow$}\fi}

\begin{document}
Optimal growth model, finite horizon\\
Main objective: to develop the tools for charactenzing denamic devisions in a model with many periods.\\
Examples: financial savings, investment in human capital, investment in technology at.\\
⇒ abstract from the static chorces le.g. consumption is. leiscure),\\
focus on the dynamic tradeoft only\\
The environment:\\
Consider a (represent ative) consumer who lives for $T+1$ periods, has access to production technology that utilizes capital, and needs to allocate the vesources betwen consumption a capital investm. in each period.

You can think of this as a Robinson-Cursol economy. OR as a decision problem of the Social Planner

Preferences: $U\left(c_{0}, c_{1}, \cdots, c_{T}\right)$ - additively separable

$$
\begin{array}{r}
=\sum_{t=0}^{T} \beta^{+} u\left(C_{4}\right), \text { where } \beta \in(0,1): \text { time discount } \\
\text { (pactor } \\
\text { (patience) }
\end{array}
$$

Typical assumptions on $U(C)$ :

\begin{itemize}
  \item $U(C)$ - cont differentiable
  \item $u^{\prime}(c)>0$ (ineceasing)
  \item $u^{n}(c)<0$ (concave)
  \item $\lim _{c \rightarrow 0} u^{\prime}(c)=+\infty \quad(\Rightarrow$ people never choose $c=0)$
\end{itemize}

\section*{Technology:}
\begin{itemize}
  \item production $F\left(k_{+}, N\right)$, assume $N=1$ (normalize) consumer $=$ moducer, decides how much be to use - assume $k_{++1} \geqslant 0, t=0, \ldots, T$
\end{itemize}

\section*{Typical assumptions on $F(k, N)$ :}
\begin{itemize}
  \item $F(k, N)$ - cont. differ
  \item $F(0, N)=0$
  \item $F_{K}(k, N)>0$ (increasing)
\end{itemize}

Sometimes we wrll assume that $F(k, 1)=A K$, e.S. The las assumption is violated

\begin{itemize}
  \item $C_{+}+i_{+} \leqslant F\left(k_{+}, N\right)$ (feasibility constraint for the conshmer)
  \item $k_{++1} \leqslant i_{t}+(1-\alpha) k_{+}$(low of motion for capital)\\
$\Rightarrow \quad C_{+}+\mu_{+}+k_{++1} \leq \mu_{+}+(1-\delta) k_{+}+F\left(k_{+}, N\right)$\\
$\Rightarrow C_{+}+k_{++1} \leqslant(1-\delta) k_{+}+F\left(k_{+}, N\right), t=0, \ldots, T$\\
feasibility constraint I resourse constraint I budget constraint
  \item $k_{0}$ is given
\end{itemize}

We will label $F(k, N)+(1-\alpha) k=f(k)$

\section*{Consumer's objestrive:}
\begin{itemize}
  \item choose the consumption and capital streams $\left(c_{+}, k_{+4}\right)_{+=0}^{\top}$
  \item in order to maximize life-time utility
  \item subject to the feasibility constraint/s) (the tradeofts imposed by the available technology)
\end{itemize}

Thus, the consumer's decision problem is:

$$
\begin{aligned}
& \max _{\left\{c_{t}, k_{t+1}\right\}_{t=0}^{T}} \sum_{t=0}^{T} \beta^{1} u\left(c_{t}\right) \\
& \begin{array}{ll}
\text { s.t. } & c_{+}+k_{++1} \leqslant f\left(k_{+}\right) \\
& k_{++1} \geqslant 0 \\
& k_{0} \text { is given }
\end{array} \quad \begin{array}{c}
\left(\lambda_{+} \beta^{+}\right) \\
\left(\mu_{+} \beta^{+}\right)
\end{array} \quad\left\{\begin{array}{l}
=\int \subset \mathbb{R}^{T+1} \\
11 \\
\left\{\begin{array}{l}
\text { (co, }, \ldots, G \\
\text { these coustr. } \\
\text { are satistical }
\end{array}\right\}
\end{array}\right.
\end{aligned}
$$

(1). The objective $f-n$ is continuast

\begin{itemize}
  \item the choice set (specified by all the constrants) is compact $\Rightarrow$ /Weierstrass's Theorem / → the max. of the objective $f-n$ is achieved in this set
\end{itemize}

$$
\begin{aligned}
\left(\exists \quad\left\{G_{4}^{*}\right\}_{t=0}^{T} \subset S: U\left(c_{0}^{*}, \ldots, G^{*}\right)=\right. & \max _{\left\{c_{4}\right\}_{t=0} \subset S} U\left(c_{0}, \ldots, G\right) \\
\text { all coustr. } & \text { all constraints }
\end{aligned}
$$

One can verify that $f(k)$-concave $\Rightarrow S$-convex\\
(2) $U\left(C_{0}, \ldots, G\right)$ - stricty concare $\Rightarrow$ the maximum is unique,\\
the solution can be obtained by the Kuhn-Tucker method\\
(3) $\lim _{c \rightarrow 0} u^{\prime}(c)=+\infty \Rightarrow c_{4}^{x}>0, t=0, \ldots, T$\\
(i.e. The solution is interior w.r.t. $G$ )

Let's fird the solution!\\
$L=\sum_{+=0}^{T} \beta^{t}\left[u\left(c_{a}\right)-\lambda_{+}\left(c_{a}+k_{++1}-f\left(k_{I}\right)\right)+\mu_{I} k_{I+1}\right]$\\
$F O C:$


\begin{align*}
& \left(c_{+}\right): \quad u^{\prime}\left(c_{+}\right)-\lambda_{+}=0, t=0 . T  \tag{1}\\
& \left(k_{++1}\right): \quad-\lambda_{+}+\mu_{+}+\beta \lambda_{++1} f^{\prime}\left(k_{++1}\right)=0, t=0, \ldots, T-1  \tag{12)}\\
& \left(k_{T+1}\right): \quad-\lambda_{T}+\mu_{T}=0 \tag{3}
\end{align*}


K- $T$ cond: $\quad \lambda_{+}\left(C_{+} k_{++1}-f\left(k_{+}\right)\right)=0, t=0, \ldots, T$

\[
\begin{array}{ll}
\mu_{+} k_{++1}=0, & t=0, \ldots, T  \tag{5}\\
\lambda_{+} \geqslant 0 & \\
\mu_{+} \geqslant 0 & , t=0, \ldots, T \\
c_{+}+k_{++1} \leqslant f\left(k_{+}\right) & \\
k_{++1} \geqslant 0 &
\end{array}
\]

Key properties of the solution:\\
(1) $k_{T+1}=0$ (no idle capital is left in the last period)

$$
\begin{aligned}
& \Gamma u^{\prime}\left(C_{T}\right)>0 \Rightarrow /(1), t=T / \Rightarrow \lambda_{T}>0 \Rightarrow \\
& /(3) / \Rightarrow \mu_{T}>0 \Rightarrow /(5), t=T / \Rightarrow k_{T+1}=0
\end{aligned}
$$

(2) $c_{+}+k_{++1} \mapsto f\left(k_{+}\right)$(the kesourse constrant holds with equality - no resourses are wasted not surprising)

$$
\begin{aligned}
& \Gamma\left(u^{\prime}\left(c_{1}\right)>0 \Rightarrow /(1) / \Rightarrow\right. \\
& \Rightarrow \lambda_{+}>0 \Rightarrow /(4) / \Rightarrow c_{+}+k_{++1}=f\left(k_{+}\right)
\end{aligned}
$$

(3) $\mu_{+}=0, t=0, \ldots, T-1$\\
$t=0, \ldots, \top$ in order to have $C_{4}>0$, it must be that $k_{+}>0 \Rightarrow$

$$
\Rightarrow \mu_{+}=0, \quad t=0, \ldots, T-1 \quad f(0)=0 \text { is key here }
$$

Thus, $(1) \Rightarrow u^{\prime}(G)=\lambda_{+}, \quad t=0, \ldots, T$

$$
\begin{aligned}
& u^{\prime}\left(c_{A+1}\right)=\lambda_{A+1}, \quad t=0, \ldots, T-1 \\
& \text { (2) } \Rightarrow \lambda_{+}=\beta \lambda_{++1} \cdot f^{\prime}\left(k_{++1}\right)+\mu_{+}, t=0, \ldots, T-1 \\
& \text { (3): } \mu_{t}=0 \\
& \Rightarrow(* \underbrace{u^{\prime}(G)=\beta u^{\prime}(G+1) \cdot f^{\prime}\left(k_{++1}\right)}_{\text {marginal utility }} \underbrace{}_{R}, t=0, \ldots, T-1 \\
& \text { of current cans. willity of } \\
& \text { future cons. } \\
& \text { future cons. can be gene- } \\
& \text { rated by giving up } 1 \\
& \text { unit of current cons. } \\
& \text { Property (2) ⇒ } C_{+}=f\left(k_{+}\right)-k_{++1} \quad \Rightarrow \\
& \text { (*) } k_{0} \text { is given }
\end{aligned}
$$

Thus $\left\{k_{+}\right\}_{+=0}^{\top}$ solves the original decision problem if and only if it solves $(* *)$.\\
(Note that $\left\{\mathrm{k}_{+}\right\}_{0}^{\top}$ pors down $\{\mathrm{G}\}_{0}^{\top}$ :

$$
G=f\left(k_{4}\right)-k_{++1}
$$

Thus we can solve for $\{G\}_{0}^{T}$ or $(k)_{0}^{T}$ only).\\
$(* *)$ is the second-order fifference equation.\\
It has a 2 -parameter fanily of solutrons (e.g. if $k_{0}, k_{1}$ are known $\Rightarrow(* *)$ pins down $k_{2}$\\
$k_{1}, k_{2}$ are known ⇒ - $\cdots$ - $k_{3}$ )\\
This solution is unique (because it is equivalent to the solution to the original problem, which has a unique solution).

But finding this solution is not trivial... I Note that we know ko $8 k_{T+1}=0$, need to find $\left.\left(k_{1}, k_{2}, \ldots, k_{T}\right\}\right)$


\end{document}